\documentclass[10pt,a4paper]{amsart}
\usepackage[margin=19mm]{geometry}
\usepackage{amsmath,amssymb,mathtools}
\usepackage[scr=rsfs,cal=euler]{mathalfa}
\usepackage{xfrac}
\usepackage[unicode,hidelinks,bookmarks=false]{hyperref}
\newcommand{\ie}{i.e.\ }
\newcommand{\cf}{cf.\ }
\newcommand{\resp}{respectively\ }
\newcommand{\Subsection}{\subsection}
\providecommand{\nobreakdash}{\nobreak\hbox{-}}
\setlength{\emergencystretch}{2em}
\multlinegap0pt
\usepackage{bm}
\usepackage{bbm}
\usepackage{dsfont}
\usepackage{xspace}
\usepackage{cancel}
\usepackage{mathtools}
\usepackage{amsthm}
\makeatletter
\@ifundefined{theorem}{\newtheorem{theorem}{Theorem}}{}
\@ifundefined{proposition}{\newtheorem{proposition}[theorem]{Proposition}}{}
\@ifundefined{restatable}{\newtheorem{restatable}{Restatable}}{}
\makeatother
\DeclarePairedDelimiter{\norm}{\|}{\|}
\newcommand{\algname}[1]{\texttt{#1}}
\title[Why Do We Need Warm-up? A Theoretical Perspective]{Why Do We Need Warm-up? A Theoretical Perspective}
\author{Foivos Alimisis, Rustem Islamov, Aurelien Lucchi}
\date{Source: arXiv:2510.03164v2 (CC BY 4.0)}
\begin{document}
\maketitle

\noindent\emph{Source and licence.} Why Do We Need Warm-up? A Theoretical Perspective. Version arXiv:2510.03164v2 (CC BY 4.0), distributed under the Creative Commons Attribution 4.0 International licence, as recorded in the arXiv licence metadata for this exact version and shown on the abstract page https://arxiv.org/abs/2510.03164v2. Full rights notice: licenses/arxiv-2510.03164-CC-BY-4.0.txt.\par\smallskip

\noindent\textit{Editorial scope.} Counted unit: the Proposition whose source label is \texttt{prop:deep\_non\_linear\_many\_nonlinearities} in the article (source0.tex lines 1565--1589, source label \texttt{prop:deep\_non\_linear\_many\_nonlinearities}), delivered together with the article's own complete original proof of it (source0.tex lines 1591--1832, one \texttt{proof} environment of 242 lines). No mathematics is written, repaired or re-derived: statement and proof are the article's own text line for line. Required context retained uncounted: prop:deep\_linear (lines 480--502); eq:Hessian\_bound\_deepLN (lines 1438--1441); eq:warmup\_constants\_linear\_network (lines 1453--1456); eq:global\_constants\_linear\_network (lines 1544--1549). Every definition, display and label of those lines is retained unchanged. Omitted from this package and not redistributed: 1--479, 503--1437, 1442--1452, 1457--1543, 1550--1564, 1590--1590, 1833--4096. Nothing inside a retained excerpt is omitted or altered: every retained body is delivered exactly as the article's lines are written, so no omission receipt of any kind accompanies this package. Citations: the retained excerpts carry no citation command at all, so no bibliography entry of the article is transcribed and no reference is redistributed. Reference adaptation: none was needed and none was made; every reference inside the delivered unit resolves inside it, so no unresolved reference or citation is produced. Printed slips kept literal: no printed word, symbol, hypothesis or number of the counted statement or of its proof was altered. Layout and class: the article's own class is \texttt{article} with packages including xcolor, amsmath, amsthm, comment, multirow, natbib, graphicx, fontenc, inputenc, threeparttable, makecell, babel, amssymb, bbm; the wrapper is the lane's pinned \texttt{amsart} editorial wrapper at 10pt on A4, which restates the theorem-like environments the retained text needs and adds the article's own macro definitions that the retained text needs (\texttt{\textbackslash algname}, \texttt{\textbackslash norm}), with extra wrapper packages (bm, bbm, dsfont, xspace, cancel, mathtools). The delivered numbering is the unit's own. Assets: no retained span contains a figure environment, a graphics-inclusion command, a PDF-page inclusion, a table or a TikZ picture; no image file is copied into this package, the package declares no asset dependency and no image file is redistributed with it. Licence: version arXiv:2510.03164v2 of this article is distributed under the Creative Commons Attribution 4.0 International licence, as recorded in the arXiv licence metadata for this exact version and shown on the abstract page https://arxiv.org/abs/2510.03164v2. A full rights notice is delivered at licenses/arxiv-2510.03164-CC-BY-4.0.txt. Characters: every retained excerpt is pure ASCII, so the delivered engine, XeLaTeX with its default Latin Modern text font, reports no missing character.
\begin{restatable}{proposition}{thmdeeplinear}
\label{prop:deep_linear}
   Consider a deep linear network with $\ell$ layers and MSE loss:
\[
    f(W) \equiv f(W_1,\dots, W_{\ell}) = \|Y- W_1 W_2 \dots W_{\ell} X\|_{\rm F}^2,
\]
where $Y \in \mathbb{R}^{c \times m}$ are the labels, $X \in \mathbb{R}^{d \times m} (d\leq m)$ is the input, and $W_i \in \mathbb{R}^{n_{i-1} \times n_i}$, where $n_0 = c$ and $n_{\ell} = d$ are layer dimensions. In the weakly balanced subspace, $\|W_1\|_F = \|W_2\|_F$, we have:

$(i)$ If $\|W_1\|_{\mathrm F}^\ell \|X\|_2 < (1-\frac{1}{\ell}) \sqrt{f(W)}$, then it holds
\begin{equation*}
\|\nabla^2 f(W)\| \leq H_0^{\text{warm-up}} + H_1^{\text{warm-up}} (f(W)-f^*)
\end{equation*}
with constants defined in Eq.~\eqref{eq:warmup_constants_linear_network}.

$(ii)$ If we additionally assume strong balancedness, then it holds
\begin{equation*}
\|\nabla^2 f(W)\| \leq H_0 + H_1 (f(W)-f^*) \hspace{2mm} \text{for all} \hspace{1mm} W,
\end{equation*}
with constants defined in Eq.~\eqref{eq:global_constants_linear_network}.



\end{restatable}

\begin{equation}
\label{eq:Hessian_bound_deepLN}
    \|\nabla^2 f(W) \|_2 \leq 2 \ell^2 \|W_1\|_{\rm F}^{2 \ell-2} \|X\|^2_2 + 2 (\ell^2-\ell) \|W_1\|_{\rm F}^{\ell-2} \|X\|_2 \sqrt{f(W)}.
\end{equation}

\begin{equation}
\label{eq:warmup_constants_linear_network}
H_0^{\text{warm-up}}:=4(\ell^2-\ell) \left(1-\frac{1}{\ell}\right)^{\frac{\ell-2}{\ell}} \|X\|_2^{\frac{2}{\ell}}\bigl(1+(f^*)^{\frac{\ell-1}{\ell}}\bigr), \qquad H_1^{\text{warm-up}}:= 4 (\ell^2-\ell) \left(1-\frac{1}{\ell}\right)^{\frac{\ell-2}{\ell}} \|X\|_2^{\frac{2}{\ell}}.
\end{equation}

\begin{align}
\label{eq:global_constants_linear_network}
 &H_0 :=A_0 \norm{X}_2^2 + A_1 \norm{X}_2 \left(\frac{f^*}{2} + \frac{1}{2}\right) + (A_2 \norm{X}_2^2+A_3 \norm{X}_2)\bigl(1+(f^*)^{\frac{\ell-1}{\ell}}\bigr), \nonumber \\
&H_1 := (A_1 + A_3) \|X\|_2 + A_2 \norm{X}_2^2
.
\end{align}

\begin{proposition}
\label{prop:deep_non_linear_many_nonlinearities}

Let $f$ be defined as
\begin{equation*}
    f(W) \equiv f(W_1,...,W_{\ell}) = \|Y- \underbrace{W_1 \phi_1(W_2 \phi(W_3 \dots \phi_{\ell-1}(W_{\ell} X) \dots ))}_F\|_F^2
\end{equation*}
where $\phi_i$ is leaky-ReLU activation function with slopes $1$ and $b_i$, i.e., $\phi_i(x)=\max \lbrace b_i x,x \rbrace$, $0 < b_i \leq 1$.
Assume that over the course of \algname{GD}:
%\begin{itemize}
    %\item
\newline $\bullet$ $\lambda_{\min}(W_i^\top W_i) \geq h_i >0,$ for $i=1,\dots,\ell-1$.
%\item 
\newline $\bullet$ The layers $W_i$ are weakly balanced, i.e., $\|W_1\|_{\rm F} = \ldots =\|W_{\ell}\|_{\rm F}.$

Then,

i) If it holds $\|W_1\|_F^{\ell} \|X\|_2 < (1-\frac{1}{\ell}) \sqrt{f(W)}$, then
we have $\|\nabla^2 f(W)\|_2 \leq H_0^{\text{warm-up}} + H_1^{\text{warm-up}} (f(W)-f^*)$, with $H_0^{\text{warm-up}}$ and $H_1^{\text{warm-up}}$ defined as in Eq.~\eqref{eq:warmup_constants_nonlinear_network}.

ii) It holds $\|\nabla^2 f(W)\|_2 \leq H_0^{\text{warm-up}} + H_1^{\text{warm-up}} (f(W)-f^*)$ for all $W$, with $H_0$ and $H_1$ defined as in Eq.~\eqref{eq:H01-post-final}.


    
\end{proposition}

\begin{proof}
The proof is divided into two parts, similarly to the proof of Proposition \ref{prop:deep_linear}: the first obtains an upper bound for the norm of the Hessian, while the second obtains a lower bound on the loss value.

The first part in the proof of Proposition \ref{prop:deep_linear} was easy, as one has ready formulas for the Hessian. In this case, the situation is more involved and we come up with a more general process to estimate the spectral norm of the Hessian based on the gradient finite differences.

\paragraph{Upper bound for the Hessian norm:}
To simplify the notation, we set
    \begin{align*}
        &Z_{\ell} = W_{\ell} X \\
        &A_{\ell-1} = \phi_{\ell-1} (Z_{\ell}) \\
        & Z_{\ell-1} = W_{\ell-1} A_{\ell-1} \\
        &\vdots 
        \\
        & Z_2 = W_2 A_2 \\ &
        A_1 = \phi_1(Z_2) \\ &
        Z_1 = W_1 A_1 = F.
    \end{align*}
    
    By the backpropagation algorithm for the gradient, we have that the gradient of $f$ can be computed as
    \begin{align*}
        \frac{\partial f}{\partial W_i} = \delta_i A_i^\top
    \end{align*}
    where $\delta_i$ is defined recursively as
    \begin{align*}
        &\delta_1 = -2(Y-F) \\ &
        \delta_2 = W_1^\top \delta_1 \odot \phi_1'(Z_2) \\ &
        \vdots \\ &
        \delta_i = W_{i-1}^\top \delta_{i-1} \odot \phi_{i-1}'(Z_i).
    \end{align*}

We need to upper bound the difference of the gradient defined in two distinct, sufficiently close points $W=(W_1,\dots,W_{\ell})$ and $\bar W = (\bar W_1,\dots, \bar W_{\ell})$. We also define
\begin{equation*}
    \textnormal{dist}(W, \bar W) := \sqrt{\sum_{i=1}^{\ell} \|W_i- \bar W_i\|^2_{\rm F}}.
\end{equation*}

It holds that
\begin{equation*}
    \|\nabla f(W)-\nabla f(\bar W)\|_{\rm F} \leq \sum_{i=1}^{\ell} \left\|\frac{\partial f}{\partial W_i}(W)-\frac{\partial f}{\partial W_i}(\bar W) \right\|_{\rm F}.
\end{equation*}

We have
\begin{equation}
\label{eq:basic_bound_partial_derivatives}
    \left\|\frac{\partial f}{\partial W_i}(W)-\frac{\partial f}{\partial W_i}(\bar W) \right\|_{\rm F} = \|\delta_i A_i^\top - \bar \delta_i \bar A_i^\top\|_{\rm F} \leq \|\delta_i\|_{\rm F} \|A_i-\bar A_i\|_{\rm F}+\|\bar A_i\|_{\rm F} \|\delta_i-\bar \delta_i\|_{\rm F}.
\end{equation}
Here we use a bar to denote the sequences of matrices related to the point $\bar W$. We deal with the four sequences appearing in this upper bound one by one, starting from $\bar A_i$. We can equivalently deal with $\bar A_i$ as the only difference will be to substitute $\bar W$ in place of $W$.

We have
\begin{align*}
    A_i =  \phi_i(W_{i+1}A_{i+1}), \hspace{1mm} \text{for} \hspace{1mm} i=1, \dots ,\ell-2,
\end{align*}
 thus
\begin{align*}
    \|A_i\|_{\rm F} =  \|\phi_i(W_{i+1}A_{i+1})\|_{\rm F} \leq \|W_{i+1} A_{i+1}\|_{\rm F} = \|W_1\|_{\rm F} \|A_{i+1}\|_{\rm F}. 
\end{align*}
The inequality follows from the fact that $\phi_i$ is leaky-ReLU, thus $|\phi_i(x)| \leq |x|$ and the last equality by the weakly balanced assumption, i.e. that $\|W_i\|_{\rm F} = \|W_1\|_{\rm F}$.

This implies that
\begin{equation}
\label{eq:bound_A_i}
    \|A_i\|_{\rm F} \leq \|W_1\|^{\ell-1-i} \|A_{\ell-1}\| = \|W_1\|^{\ell-1-i} \|\phi_{\ell-1}(W_{\ell} X)\|_{\rm F} \leq \|W_1\|_{\rm F}^{\ell-i} \|X\|_2.
\end{equation}

Similarly, it holds
\begin{equation}
\label{eq:bar_A_i}
    \|\bar A_i\|_{\rm F} \leq \|\bar W_1\|_{\rm F}^{\ell-i} \|X\|_2.
\end{equation}

Now, we deal with $A_i-\bar A_i$:
\begin{align*}
    &\|A_i-\bar A_i\|_{\rm F} = \nonumber\|\phi_i(W_{i+1}A_{i+1}) - \phi_i(\bar W_{i+1} \bar A_{i+1}) \|_{\rm F} \leq \|W_{i+1}A_{i+1} - \bar W_{i+1} \bar A_{i+1}\|_{\rm F} \leq \\ & \|A_{i+1}\|_{\rm F} \|W_{i+1}-\bar W_{i+1}\|_{\rm F} + \|\bar W_{i+1}\|_{\rm F} \|A_{i+1} - \bar A_{i+1}\|_{\rm F}  \leq \\ & \|A_{i+1}\|_{\rm F} \textnormal{dist}(W,\bar W) + \|\bar W_1\|_{\rm F} \|A_{i+1}-\bar A_{i+1}\|_{\rm F}.
\end{align*}

By an induction argument, we can get the bound
\begin{equation*}
    \|A_i-\bar A_i\|_{\rm F} \leq \left(\sum_{k=i+1}^{\ell-1} \|A_k\| \|\bar W_1\|^{k-i-1}\right) \textnormal{dist}(W,\bar W) + \|\bar W_1\|_{\rm F}^{\ell-i-1} \|A_{\ell-1}-\bar A_{\ell-1} \|_{\rm F}
\end{equation*}
and by inequality \eqref{eq:bound_A_i}, we have
\begin{align}
\label{eq:A_i-barA_i}
     \|A_i-\bar A_i\|_{\rm F} & \leq \nonumber \left(\sum_{k=i+1}^{\ell-1} \|W_1\|_{\rm F}^{\ell-k} \|\bar W_1\|^{k-i-1}\right) \textnormal{dist}(W,\bar W) \|X\|_2 + \|\bar W_1\|_{\rm F}^{\ell-i-1} \|W_{\ell}-\bar W_{\ell}\|_{\rm F} \|X\|_2 \\ & \leq \left(\sum_{k=i+1}^{\ell-1} \|W_1\|_{\rm F}^{\ell-k} \|\bar W_1\|^{k-i-1}\right) \textnormal{dist}(W,\bar W) \|X\|_2 + \|\bar W_1\|_{\rm F}^{\ell-i-1} \textnormal{dist}(W,\bar W) \|X\|_2.
\end{align}

Now we move to $\delta_i$. It holds
\begin{equation*}
    \|\delta_i\|_{\rm F} = \|W_{i-1}^\top \delta_{i-1} \odot \phi_{i-1}'(Z_i)\|_{\rm F} \leq \|W_{i-1}\|_{\rm F} \|\delta_{i-1}\|_{\rm F} = \|W_1\|_{\rm F} \|\delta_{i-1}\|_{\rm F}.
\end{equation*}
This implies that
\begin{equation*}
     \|\delta_i\|_{\rm F} \leq \|W_1\|_{\rm F}^{i-1} \|\delta_1\|_{\rm F} = 2 \|W_1\|_{\rm F}^{i-1} \sqrt{f(W)}
\end{equation*}
and similarly
\begin{equation}
\label{eq:bar_delta}
    \|\bar \delta_i\|_{\rm F} \leq 2 \|\bar W_1\|_{\rm F}^{i-1} \sqrt{f(\bar W)}.
\end{equation}

For the sequence $\delta_i-\bar \delta_i$, we have
\begin{align*}
    \|\delta_i - \bar \delta_i\|_{\rm F} = \| W_{i-1}^\top \delta_{i-1} \odot \phi_{i-1}'(Z_i) - \bar W_{i-1}^\top \bar \delta_{i-1} \odot \phi_{i-1}'(\bar Z_i)\|_{\rm F}
\end{align*}
and since all entries of $Z_i$ are non-zero and $\bar Z_i$ is taken sufficiently close to $Z_i$, these two points feature the same activation pattern, thus $\phi_{i-1}'(Z_i) = \phi_{i-1}'(\bar Z_i)$. This gives
\begin{align*}
    \|\delta_i - \bar \delta_i\|_{\rm F} & \leq \| W_{i-1}^\top \delta_{i-1} - \bar W_{i-1}^\top \bar \delta_{i-1}\|_{\rm F} \leq \|W_{i-1}\|_{\rm F} \| \delta_{i-1}-\bar \delta_{i-1}\|_{\rm F} + \|\bar \delta_{i-1}\|_{\rm F} \|W_{i+1}-\bar W_{i+1}\|_{\rm F} 
 \\ &\leq \|W_1\|_{\rm F} \| \delta_{i-1}-\bar \delta_{i-1}\|_{\rm F} + \|\bar \delta_{i-1}\|_{\rm F} \textnormal{dist}(W,\bar W).
\end{align*}

By induction, we have
\begin{align*}
    \|\delta_i - \bar \delta_i\|_{\rm F} & \leq \sum_{k=i-1}^1 \|\bar \delta_k\|_{\rm F} \| W_1 \|_{\rm F}^{i-1-k} \textnormal{dist}(W,\bar W) + \|W_1\|_{\rm F}^{i-1} \| \delta_1 - \bar \delta_1\|_{\rm F} \\ & \leq 2 \sqrt{f(\bar W)} \sum_{k=i-1}^1  \|\bar W_1\|_{\rm F}^{k-1} \| W_1 \|_{\rm F}^{i-1-k} \textnormal{dist}(W,\bar W) + \|W_1\|_{\rm F}^{i-1} \| \delta_1 - \bar \delta_1\|_{\rm F}. 
\end{align*}
The second inequality in the previous derivation follows by inequality \eqref{eq:bar_delta}.

For $\|\delta_1-\bar \delta_1\|_{\rm F}$, we have
\begin{align*}
  &\|\delta_1-\bar \delta_1\|_{\rm F} = 2 \|W_1 A_1-\bar W_1 \bar A_1\|_{\rm F} \leq 2\|W_1\|_{\rm F}\|A_1-\bar A_1\|_{\rm F}+ 2 \|\bar A_1\|_{\rm F} \|W_1-\bar W_1\|_{\rm F} \leq \\ &  2\|W_1\|_{\rm F} \left(\left(\sum_{k=2}^{\ell-1} \|W_1\|_{\rm F}^{\ell-k} \|\bar W_1\|^{k-2}\right) + \|\bar W_1\|_{\rm F}^{\ell-2} \right) \textnormal{dist}(W,\bar W) \|X\|_2 + 2\|\bar W_1\|_{\rm F}^{\ell-1} \textnormal{dist}(W,\bar W) \|X\|_2 = \\ & 2\left(\|W_1\|_{\rm F} \left(\left(\sum_{k=2}^{\ell-1} \|W_1\|_{\rm F}^{\ell-k} \|\bar W_1\|^{k-2}\right) + \|\bar W_1\|_{\rm F}^{\ell-2} \right) + \|\bar W_1\|_{\rm F}^{\ell-1} \right) \textnormal{dist}(W,\bar W) \|X\|_2.
\end{align*}

Thus,
\begin{align}
\label{eq:eq8}
    &\|\delta_i - \bar \delta_i\|_{\rm F}  \nonumber \leq 2 \sqrt{f(\bar W)} \sum_{k=i-1}^1  \|\bar W_1\|_{\rm F}^{k-1} \| W_1 \|_{\rm F}^{i-1-k} \textnormal{dist}(W,\bar W) + \\& 2 \|W_1\|_{\rm F}^{i-1} \left(\|W_1\|_{\rm F} \left(\left(\sum_{k=2}^{\ell-1} \|W_1\|_{\rm F}^{\ell-k} \|\bar W_1\|^{k-2}\right) + \|\bar W_1\|_{\rm F}^{\ell-2} \right) + \|\bar W_1\|_{\rm F}^{\ell-1} \right) \textnormal{dist}(W,\bar W) \|X\|_2. 
\end{align}

Combining inequalities \eqref{eq:basic_bound_partial_derivatives},\eqref{eq:bar_A_i},\eqref{eq:A_i-barA_i},\eqref{eq:bar_delta} and \eqref{eq:eq8},
we get
\begin{align*}
    &\left\|\frac{\partial f}{\partial W_i}(W)-\frac{\partial f}{\partial W_i}(\bar W) \right\|_{\rm F} \leq \\ &  2\|\bar W_1\|_{\rm F}^{i-1} \sqrt{f(\bar W)} \left(\left(\sum_{k=i+1}^{\ell-1} \|W_1\|_{\rm F}^{\ell-k} \|\bar W_1\|^{k-i-1}\right) + \|\bar W_1\|_{\rm F}^{\ell-i-1}\right)  \textnormal{dist}(W,\bar W) \|X\|_2 + \\ & 2 \|\bar W_1\|_{\rm F}^{\ell-i} \|X\|_2 \sqrt{f(\bar W)} \sum_{k=i-1}^1  \|\bar W_1\|_{\rm F}^{k-1} \| W_1 \|_{\rm F}^{i-1-k} \textnormal{dist}(W,\bar W) + \\ & 2\|\bar W_1\|_{\rm F}^{\ell-i} \|W_1\|_{\rm F}^{i-1} \left(\|W_1\|_{\rm F} \left(\left(\sum_{k=2}^{\ell-1} \|W_1\|_{\rm F}^{\ell-k} \|\bar W_1\|^{k-2}\right) + \|\bar W_1\|_{\rm F}^{\ell-2} \right) + \|\bar W_1\|_{\rm F}^{\ell-1} \right) \textnormal{dist}(W,\bar W) \|X\|^2_2,  
\end{align*}
thus
\begin{align*}
    &\frac{\left\|\frac{\partial f}{\partial W_i}(W)-\frac{\partial f}{\partial W_i}(\bar W) \right\|_{\rm F}}{\textnormal{dist}(W,\bar W)} \leq \\ &  2\|\bar W_1\|_{\rm F}^{i-1} \sqrt{f(\bar W)} \left(\left(\sum_{k=i+1}^{\ell-1} \|W_1\|_{\rm F}^{\ell-k} \|\bar W_1\|^{k-i-1}\right) + \|\bar W_1\|_{\rm F}^{\ell-i-1}\right) \|X\|_2 + \\ & 2 \|\bar W_1\|_{\rm F}^{\ell-i} \|X\|_2 \sqrt{f(\bar W)} \sum_{k=i-1}^1  \|\bar W_1\|_{\rm F}^{k-1} \| W_1 \|_{\rm F}^{i-1-k} + \\ & 2 \|\bar W_1\|_{\rm F}^{\ell-i} \|W_1\|_{\rm F}^{i-1} \left(\|W_1\|_{\rm F} \left(\left(\sum_{k=2}^{\ell-1} \|W_1\|_{\rm F}^{\ell-k} \|\bar W_1\|^{k-2}\right) + \|\bar W_1\|_{\rm F}^{\ell-2} \right) + \|\bar W_1\|_{\rm F}^{\ell-1} \right) \|X\|^2_2
\end{align*}
and taking the limit as $\bar W \longrightarrow W$, we get
\begin{align*}
    &\lim_{\bar W \rightarrow W}\frac{\left\|\frac{\partial f}{\partial W_i}(W)-\frac{\partial f}{\partial W_i}(\bar W) \right\|_{\rm F}}{\textnormal{dist}(W,\bar W)} \leq \\ &  2 (\ell-i) \|X\|_2 \|W_1\|_{\rm F}^{\ell-2} \sqrt{f(W)}  +2 (i-1) \|X\|_2 \|W_1\|_{\rm F}^{\ell-2} \sqrt{f(W)} + 2 (\ell-1) \|W_1\|_{\rm F}^{2 \ell-2} \|X\|_2^2 = \\ & 2 (\ell-1) \|X\|_2 \sqrt{f(W)} \|W_1\|_{\rm F}^{\ell-2} + 2 \ell \|W_1\|_{\rm F}^{2 \ell-2} \|X\|_2^2.
\end{align*}
This is because, when $\bar W \longrightarrow W$, it holds $\bar W_1\longrightarrow W_1$.

For the total gradient difference, we have
\begin{align*}
  \lim_{\bar W \rightarrow W}\frac{\left\|\nabla f(W)-\nabla f(\bar W) \right\|_{\rm F}}{\textnormal{dist}(W,\bar W)} & \leq \sum_{i=1}^\ell \lim_{\bar W \rightarrow W}\frac{\left\|\frac{\partial f}{\partial W_i}(W)-\frac{\partial f}{\partial W_i}(\bar W) \right\|_{\rm F}}{\textnormal{dist}(W,\bar W)} \\ &  \leq 2 \ell (\ell-1) \|W_1\|_{\rm F}^{\ell-2} \|X\|_2 \sqrt{f(W)} + 2 \ell^2 \|W_1\|_{\rm F}^{2 \ell-2} \|X\|_2^2.
\end{align*}
It holds
\begin{equation*}
    \|\nabla^2 f(W) \|_2 = \lim_{\bar W \rightarrow W}\frac{\left\|\nabla f(W)-\nabla f(\bar W) \right\|_{\rm F}}{\textnormal{dist}(W,\bar W)},
\end{equation*}
thus
\begin{equation}
\label{eq:Hessian_bound_non_linear}
    \|\nabla^2 f(W)\|_2 \leq 2 \ell (\ell-1) \|W_1\|_{\rm F}^{\ell-2} \|X\|_2 \sqrt{f(W)} + 2 \ell^2 \|W_1\|_{\rm F}^{2 \ell-2} \|X\|_2^2.
\end{equation}

Notice that this is the same upper bound as the one provided in \eqref{eq:Hessian_bound_deepLN}.

Thus, if $\|W\|_F^{\ell} \|X\|_2 \leq (1-\frac{1}{\ell}) \sqrt{f(W)}$, then we get $\|\nabla^2 f(W)\| \leq H_0^{\text{warm-up}} + H_1^{\text{warm-up}} (f(W)-f^*)$ with exactly the same constants as in the previous proposition, i.e.
\begin{equation}
\label{eq:warmup_constants_nonlinear_network}
H_0^{\text{warm-up}}:=4(\ell^2-\ell) \left(1-\frac{1}{\ell}\right)^{\frac{\ell-2}{\ell}} \|X\|_2^{\frac{2}{\ell}}\bigl(1+(f^*)^{\frac{\ell-1}{\ell}}\bigr), \qquad H_1^{\text{warm-up}}:= 4 (\ell^2-\ell) \left(1-\frac{1}{\ell}\right)^{\frac{\ell-2}{\ell}} \|X\|_2^{\frac{2}{\ell}}.
\end{equation}






We now move to a global bound and for that we need a global lower bound for the loss value. For $i=1 ,\dots, \ell-2$, we have
\begin{align*}
     \|W_i A_i\|_F^2 \geq \lambda_{\min}(W_i^T W_i) \|A_i\|_F^2 \geq h_i \|\phi_i(W_{i+1} A_{i+1})\|^2_F \geq h_i b_i^2 \|W_{i+1} A_{i+1}\|_F^2
\end{align*}
and by induction,
\begin{align}
\label{eq:deep_leaky_relu_basic}
    \|W_1 A_1\|_F^2 & \geq \nonumber\left(\Pi_{i=1}^{\ell-2} h_i b_i^2 \right) \|W_{\ell-1} A_{\ell-1}\|_F^2 \\ & \geq \nonumber\left(\Pi_{i=1}^{\ell-2} h_i b_i^2 \right) \lambda_{\min}(W_{\ell-1}^T W_{\ell-1}) \|A_{\ell-1}\|_F^2 \\ & = \nonumber \left(\Pi_{i=1}^{\ell-2} h_i b_i^2 \right) \lambda_{\min}(W_{\ell-1}^T W_{\ell-1}) \|\phi_{\ell-1}(W_{\ell} X)\|_F^2 \\ & \geq \nonumber \left(\Pi_{i=1}^{\ell-2} h_i b_i^2 \right) h_{\ell-1} b_{\ell-1}^2 \lambda_{\min}(X X^T) \|W_{\ell}\|_F^2 \\ & = \nonumber\left(\Pi_{i=1}^{\ell-2} h_i b_i^2 \right) h_{\ell-1} b_{\ell-1}^2 \lambda_{\min}(X X^T) \|W_1\|_F^2 \\ & = \left(\Pi_{i=1}^{\ell-1} h_i b_i^2 \right) \lambda_{\min}(X X^T) \|W_1\|_F^2.
\end{align}
We have repeatedly used the assumption that $\lambda_{\min}(W_i W_i^T) \geq h_i$ and that
\begin{equation*}
    \|\phi_i(S)\|^2_F \geq b_i^2 \|S\|_F^2,
\end{equation*}
for any matrix $S$.

To derive inequality \eqref{eq:deep_leaky_relu_basic}, we also used the weak balancedness assumption, that is, all $\|W_i\|_F$ are equal.

Defining \begin{equation}
\label{eq:K_gd_def}
K_{gd} := \left( \lambda_{\min}(XX^\top) \prod_{i=1}^{\ell-1} h_i b_i^2 \right)^{-\frac{1}{2}}
\end{equation}
we get
$$ \|W_1\|_F \leq K_{gd} \|W_1 A_1\|_F. $$

By the triangle inequality, $\norm{W_1 A_1}_F \le \norm{Y}_F + \sqrt{f(W)}$. Thus:
\begin{equation}
\label{eq:W1_bound_nonlinear}
\norm{W_1}_F \le K_{gd} \left( \norm{Y}_F + \sqrt{f(W)} \right).
\end{equation}

Now we substitute this into the sharp Hessian bound \eqref{eq:Hessian_bound_non_linear}.
\paragraph{Term 1:}
\begin{align*}
2\ell^2 \norm{W_1}_F^{2\ell-2} \norm{X}_2^2
\le 2\ell^2 K_{gd}^{2\ell-2} \left( \norm{Y}_F + \sqrt{f} \right)^{2\ell-2} \norm{X}_2^2 
\le 2\ell^2 K_{gd}^{2\ell-2} 2^{2\ell-3} \left( \norm{Y}_F^{2\ell-2} + f^{\ell-1} \right) \norm{X}_2^2.
\end{align*}
(Using Jensen's inequality $(a+b)^p \le 2^{p-1}(a^p+b^p)$ with $p=2\ell-2$).

\paragraph{Term 2:}
\begin{align*}
2(\ell^2-\ell) \norm{W_1}_F^{\ell-2} \norm{X}_2 \sqrt{f}
&\le 2(\ell^2-\ell) K_{gd}^{\ell-2} \left( \norm{Y}_F + \sqrt{f} \right)^{\ell-2} \norm{X}_2 \sqrt{f} \\
&\le 2(\ell^2-\ell) K_{gd}^{\ell-2} 2^{\ell-3} \left( \norm{Y}_F^{\ell-2} + f^{\frac{\ell-2}{2}} \right) \norm{X}_2 \sqrt{f} \\
&= 2(\ell^2-\ell) K_{gd}^{\ell-2} 2^{\ell-3} \left( \norm{Y}_F^{\ell-2}\sqrt{f} + f^{\frac{\ell-1}{2}} \right) \norm{X}_2.
\end{align*}


We now use the inequality $u \leq 1+ u^p$, which holds for all $u \geq 0$ and $p \geq 1$. In our case, it translates to 
\begin{equation*}
    \sqrt{f} \leq 1+f^{\ell-1}
\end{equation*}
and
\begin{equation*}
    f^{\frac{\ell-1}{2}} \leq 1 + f^{\ell-1}.
\end{equation*}

Term 2 can be bounded as
\begin{align*}
    & 2(\ell^2-\ell) K_{gd}^{\ell-2} 2^{\ell-3} \left( \norm{Y}_F^{\ell-2}\sqrt{f} + f^{\frac{\ell-1}{2}} \right) \norm{X}_2 \leq 2(\ell^2-\ell) K_{gd}^{\ell-2} 2^{\ell-3} \left(1+\|Y\|_F^{\ell-2} \right) \|X\|_2 \left(1 + f^{\ell-1} \right) = \\ & 2(\ell^2-\ell) K_{gd}^{\ell-2} 2^{\ell-3} \left(1+\|Y\|_F^{\ell-2} \right) \|X\|_2 + 2(\ell^2-\ell) K_{gd}^{\ell-2} 2^{\ell-3} \left(1+\|Y\|_F^{\ell-2} \right) \|X\|_2 f^{\ell-1}.
\end{align*}

Summing with the bound of Term 1 yields the inequality $\|\nabla^2 f(W)\|_2 \leq H_0 + H_1 (f(W)-f^*)$, with
\begin{align}
\label{eq:H01-post-final}
H_0 &:= 2\ell^2 K_{gd}^{2\ell-2} 2^{2\ell-3} \norm{Y}_F^{2\ell-2} \norm{X}_2^2 + 2(\ell^2-\ell) K_{gd}^{\ell-2} 2^{\ell-3} (1+ \norm{Y}_F^{\ell-2}) \norm{X}_2, \nonumber \\
H_1 &:= 2\ell^2 K_{gd}^{2\ell-2} 2^{2\ell-3} \norm{X}_2^2 + 2(\ell^2-\ell) K_{gd}^{\ell-2} 2^{\ell-3} (1+ \|Y\|_F^{\ell-2}) \norm{X}_2,
\end{align}
where $K_{gd}$ is defined as in Eq.~\eqref{eq:K_gd_def}.




    
\end{proof}

\end{document}
